%
% teil3.tex -- Beispiel-File für Teil 5
%
% (c) 2020 Prof Dr Andreas Müller, Hochschule Rapperswil
%
% !TEX root = ../../buch.tex
% !TEX encoding = UTF-8
%
\section{Matrixdarstellung
\label{moebius:section:teil5}}
\kopfrechts{Matrixdarstellung}

\subsection{Einfachere Darstellung von Möbius-Transformationen}
Möbius-Transformationen können oft als Verkettung mehrerer einfacher Transformationen
(z.\,B. Skalierungen, Drehungen, Translationen und Inversionen) dargestellt werden.
Diese direkte Schreibweise kann jedoch schnell unübersichtlich werden.
Eine deutlich einfachere und strukturiertere Darstellung erhält man durch die Verwendung von Matrizen.
\index{Matrix!einer Möbiustransformation}%
Jede Möbius-Transformation
\[
f(z) = \frac{az + b}{cz + d}
\]
kann eindeutig (bis auf Skalierung) durch die Matrix
\[
M = \begin{pmatrix} a & b \\ c & d \end{pmatrix}
\]
beschrieben werden.
Dabei gilt die Bedingung
\[
ad - bc \neq 0,
\]
da die Determinante der Matrix ungleich Null sein muss. Nur dann ist die Matrix invertierbar und die Transformation wohldefiniert.
\index{Determinante}%

\subsection{Verkettung von Transformationen}

Werden mehrere Möbius-Transformationen hintereinander ausgeführt, z.\,B.
\[
f_3(f_2(f_1(z))),
\]
so entspricht dies der Matrixmultiplikation:
\[
M = M_3 \cdot M_2 \cdot M_1.
\]

\begin{beispiel}
Gegeben seien die Transformationen:
\begin{align*}
f_1(z) &= 2e^{i\pi/4} \cdot z
\quad \Rightarrow \quad
M_1 = \begin{pmatrix} 2e^{i\pi/4} & 0 \\ 0 & 1 \end{pmatrix} 	
\\
f_2(z) &= z + (3 - 2i)
\quad \Rightarrow \quad
M_2 = \begin{pmatrix} 1 & 3 - 2i \\ 0 & 1 \end{pmatrix} 
\\
f_3(z) &= \frac{z + 1}{2z + 3}
\quad \Rightarrow \quad
M_3 = \begin{pmatrix} 1 & 1 \\ 2 & 3 \end{pmatrix}
\end{align*}
Zuerst berechnen wir:
\[
M_{21} = M_2 \cdot M_1
= \begin{pmatrix} 1 & 3 - 2i \\ 0 & 1 \end{pmatrix}
\cdot
\begin{pmatrix} 2e^{i\pi/4} & 0 \\ 0 & 1 \end{pmatrix}
=
\begin{pmatrix}
	2e^{i\pi/4} & 3 - 2i \\
	0 & 1
\end{pmatrix}
\]
\[
M = M_3 \cdot M_{21}
=
\begin{pmatrix} 1 & 1 \\ 2 & 3 \end{pmatrix}
\cdot
\begin{pmatrix}
	2e^{i\pi/4} & 3 - 2i \\
	0 & 1
\end{pmatrix}
\]
Wie berechnen die Einträge, beginnend oben Links:
\[
1 \cdot 2e^{i\pi/4} + 1 \cdot 0 = 2e^{i\pi/4}.
\]
Oben rechts:
\[
1 \cdot (3 - 2i) + 1 \cdot 1 = 4 - 2i.
\]
Unten links:
\[
2 \cdot 2e^{i\pi/4} + 3 \cdot 0 = 4e^{i\pi/4}.
\]
Unten rechts:
\[
2 \cdot (3 - 2i) + 3 \cdot 1 = 9 - 4i.
\]
Damit ergibt sich die Gesamtmatrix
\[
M =
\begin{pmatrix}
	2e^{i\pi/4} & 4 - 2i \\
	4e^{i\pi/4} & 9 - 4i
\end{pmatrix}.
\]
Diese Matrix kann nun wieder in die allgemeine Form
\[
f(z) = \frac{az + b}{cz + d}
\]
 der Möbius-Transformation eingesetzt werden, daraus ergibt sich die zugehörige Transformation
\[
f(z) =
\frac{2e^{i\pi/4} \cdot z + (4 - 2i)}
{4e^{i\pi/4} \cdot z + (9 - 4i)}.\qedhere
\]
\end{beispiel}
